\noindent\textbf{Elección del procedimiento.} Separar componentes: $\ddot x=0$, $\ddot y=-g$. Con origen en el lanzamiento se obtiene $x=v_0\cos\theta\,t$ e $y=v_0\sin\theta\,t-gt^2/2$. Imponer $y=0$ da dos raíces: $t=0$ y el regreso al suelo; se elige la segunda positiva.

\begin{enumerate}

\item La componente horizontal permanece constante en el vacío.

\[\fitmath{\mathrm{vx}=v_{0}\cdot \cos\left(\theta\right)}\]

\[\fitmath{\mathrm{vx}=(30)\cdot \cos\left((0.6981317008)\right)\approx 22.98133\,\mathrm{m/s}}\]

\item La componente vertical inicial fija el tiempo de vuelo.

\[\fitmath{\mathrm{vy}=v_{0}\cdot \sin\left(\theta\right)}\]

\[\fitmath{\mathrm{vy}=(30)\cdot \sin\left((0.6981317008)\right)\approx 19.28363\,\mathrm{m/s}}\]

\item Dividir $t(v_{y0}-gt/2)=0$ por la raíz inicial ya identificada.

\[\fitmath{\mathrm{tf}=\frac{2\cdot \mathrm{vy}}{g}}\]

\[\fitmath{\mathrm{tf}=\frac{2\cdot (19.28362829)}{(9.81)}\approx 3.931423\,\mathrm{s}}\]

\item Introducir ese tiempo en $x(t)$; equivale a $v_0^2\sin(2\theta)/g$.

\[\fitmath{R=\mathrm{vx}\cdot \mathrm{tf}}\]

\[\fitmath{R=(22.98133329)\cdot (3.931422689)\approx 90.34934\,\mathrm{m}}\]

\end{enumerate}
